% RQ4. Nguồn: results/apple-m5/barrier.csv (trung vị 7 lần), bảng barrier, sleeplatency.
\subsubsection{RQ4: Chi phí một barrier episode}

\begin{rqanswer}\textbf{Trả lời RQ4.} Với chờ spin, một episode tốn thêm khoảng 53~ns cho mỗi thread
trong cụm 4 P-core, dốc gần gấp 6 khi thread tràn sang E-core, và sụp khoảng 2\,000 lần khi số thread
vượt số core. Khi mỗi thread có core riêng, spin thuần nhanh nhất. Khi vượt số core,
\PCname{sr_bo_yield} nhanh nhất. Kết luận chỉ áp cho \PCname{apple-m5}.\end{rqanswer}

Nhóm đo số episode (một lần cả $p$ thread cùng qua barrier) mỗi giây khi giữa hai barrier không có việc
gì (Hình~\ref{fig:apple-m5:barrier_rate}, Bảng~\ref{tab:apple-m5:barrier}).

\PCfigure{apple-m5}{barrier_rate}
\PCtable{apple-m5}{barrier}

Với \PCname{sr}, mỗi thread thêm vào tốn khoảng 53~ns từ 2 tới 4 thread và khoảng 295~ns từ 4 tới
10 thread.
% nguồn: 1e9/trung vị ops_per_s của sr: p=2 94,5; p=4 200,1; p=10 1969 ns.
Phần tuyến tính đến từ $p$ phép \texttt{fetch\_add} lần lượt trên cache line của bộ đếm và $p-1$ lần
đọc lại cờ sau khi cờ đảo. Chỗ gãy sau 4 thread trùng với số P-core, và kết quả lock gãy ở cùng chỗ
ngay cả ở ô đối chứng mà lock gần như rảnh (RQ2). Vì vậy nó đến từ máy chứ không từ thuật toán barrier.
Nhóm giả thuyết hai nguyên nhân là E-core chậm hơn và cache line phải đi giữa hai cụm L2. Vì không gắn
được thread vào core, nhóm không tách được hai nguyên nhân này.

Từ 10 lên 12 thread, \PCname{sr} chậm đi khoảng 2\,070 lần, và mỗi thread vượt số core thêm khoảng
2~ms vào một episode.
% nguồn: sr: tốc độ p=10/p=12 = 2067; ms mỗi episode/(p-10), p=12..20: 1,7-2,0.
Một episode cần \emph{mọi} thread được chạy, nhưng thread đã tới quay giữ core tới hết lượt lập lịch.
Barrier giống MCS ở điểm này (RQ2).

Chờ chặn hoà vốn đúng tại số core. Dưới đó, \PCname{sr_block} luôn chậm hơn \PCname{sr}, vì đánh thức
qua kernel tốn vài micro giây. Ở 12 thread, nó nhanh hơn khoảng 217 lần.
% nguồn: trung vị ops_per_s p=12: sr_block 53319, sr 245,7.
\PCname{sr_bo_yield} nhanh hơn \PCname{std_barrier} ở mọi số thread từ 2 tới 20, và nhanh nhất trong
mọi barrier khi vượt số core (11,6~$\mu$s mỗi episode ở 12 thread). Thread chờ gọi
\texttt{sched\_yield}, nhường core cho đúng thread chưa tới barrier. Nhờ vậy episode xong mà không phải chờ hết
lượt lập lịch, và không phải trả một lần đánh thức qua kernel như \PCname{sr_block}. Khi mỗi thread có
core riêng, lời gọi hệ thống này chỉ làm thread chậm nhận ra cờ đã đảo, nên \PCname{sr_bo_yield} chậm hơn
\PCname{sr} tới 1,7 lần (8 thread).
% nguồn: ns/episode (1e9/trung vị ops_per_s): p=12 sr_bo_yield 11555, std_barrier 13657, sr_block 18755; p=8 sr_bo_yield 1225, sr 733.

Backoff làm barrier chậm đi khi mỗi thread có core riêng (\PCname{sr_bo_spin} chậm hơn \PCname{sr}
2,3 tới 9,5 lần), vì thread chờ chỉ đọc cờ nên không có tranh chấp để bớt. Một episode của
\PCname{sr_bo_sleep} ở 2 thread tốn 2,49~$\mu$s, xấp xỉ thời gian ngủ thật ngắn nhất của
\texttt{sleep\_for} khi yêu cầu 100~ns (2,54~$\mu$s).
% nguồn: sr_bo_sleep p=2: 1e9/401186 = 2493 ns; sleep_for 100 ns: nhỏ nhất 2540 ns.

Đối chiếu với dự đoán, D4.1 chỉ đúng từng đoạn, vì chi phí tăng tuyến tính theo $p$ nhưng độ dốc đổi
ở 4 thread. D4.2 đúng. Spin sụp ngay khi vượt số core, và chờ chặn chỉ thắng ở vùng đó. D4.3 đúng.
Backoff làm barrier chậm khi mỗi thread có core riêng. Các ô sụp đổ dựa trên 50--69 episode mỗi lần chạy nên chỉ cho bậc độ lớn. Khoảng
cách 2--3 bậc vẫn lớn hơn nhiều độ phân tán giữa 7 lần lặp.
% nguồn: cột iters của sr, p=12..20.
